\begin{frame}{세 갈래를 가르는 ``데이터의 형상'' --- 표로 보면 된다}
  \begin{itemize}
    \item \kb{지도}: \textbf{2개 블록} --- 왼쪽 $x$ 블록(특징 열),
      오른쪽 $y$ 블록(라벨 1열). 행 $m$ 개, 열 $(p+1)$ 개. 핵심 질문:
      ``라벨 열이 있는가?''
    \item \kb{비지도}: \textbf{1개 블록} --- $x$ 블록만 있고 라벨
      \textbf{열이 없다}. 행 $m$ 개, 열 $p$ 개. 목표: ``구조를 찾자''.
    \item \kb{강화}: 형상 자체가 다름 --- 각 행이
      $(s_t, a_t, r_t, s_{t+1})$ 4개 필드(상태, 행동, \textbf{보상},
      다음 상태)를 가진 \textbf{경험 튜플} 목록. 라벨 열 대신
      \textbf{보상 열}이 있다.
  \end{itemize}
  \vskip0.5em
  \begin{block}{``라벨 열''과 ``보상 열''의 차이}
    라벨 = ``이 입력 $x$ 에 대한 \textbf{정답}''. \; 보상 = ``이
    행동을 취한 \textbf{결과}가 좋았는가''를 숫자로 알려줄 뿐,
    ``정답 행동은 무엇이었는가''를 \textbf{아주 느리고 부분적으로}만
    알려준다. 이것이 세 갈래의 가장 근본적인 차이 (ML2에서 본격 다룸).
  \end{block}
\end{frame}
